\uuid{KrXd}
\titre{Primitive}
\auteur{bodin ; croizat ; sacre}
\organisation{exo7}
\licence{CC BY-NC-SA 4.0}
\datecreate{2022-03-08}
\niveau{L1}
\module{Analyse}
\chapitre{Calcul d'intégrales}
\sousChapitre{Primitives diverses}
\difficulte{1}
\theme{Primitive}
\exo7qcmid{hilisit.equations-differentielles.001}

\contenu{
\begin{question}
\qclassification{127.11}
Quelles sont les affirmations vraies ?
\begin{answers}
  \bad{$x^3$ est une primitive de $3x^2+3$.}
  \good{$x^3+3$ est une primitive de $3x^2$.}
  \bad{ $\ln(x^2+1)$ est une primitive de $\frac 1{x^2+1}$.}
  \good{$\sqrt x$ est une primitive de $\frac 1{2\sqrt x}$ (sur $]0,+\infty[$).}
\end{answers}
\begin{explanations}
Pour vérifier si une fonction \(f\) est une primitive d'une fonction \(g\), on calcule la dérivée de  \(f\) et on regarde si on obtient bien la fonction \(g\). La dérivée de $x^3$ et de $x^3+3$ est $3x^2$. La dérivée de $\ln(x^2+1)$ est $\frac {2x}{x^2+1}$ et non $\frac 1{x^2+1}$. La dérivée de $\sqrt x$ sur $]0,+\infty[$ est bien $\frac 1{2\sqrt x}$.
\end{explanations}
\end{question}
}
